\documentclass[10pt]{beamer}

% --- Modern Theme ---
\usetheme{metropolis}
\usecolortheme{default}
\setbeamercolor{background canvas}{bg=white}

% --- Packages ---
\usepackage[utf8]{inputenc}
\usepackage{booktabs}
\usepackage{array}
\usepackage{tabularx}
\usepackage{amsmath}
\usepackage{amssymb}
\usepackage[scale=2]{ccicons}
\usepackage{pgfplots}
\usepgfplotslibrary{dateplot}
\usepackage{xspace}

% --- Title Information ---
\title{Foundations of Artificial Intelligence}
\subtitle{Part VI: Introduction to Neural Networks}
\institute{Department of Computer Science}
\date{}

\begin{document}

\maketitle

\begin{frame}{Lecture Roadmap}
    \setbeamertemplate{section in toc}[sections numbered]
    \tableofcontents[hideallsubsections]
\end{frame}

\begin{frame}{Learning Objectives}
    By the end of this lecture, you should be able to:
    \begin{itemize}
        \item Describe the biological inspiration behind artificial neural networks.
        \item Define a perceptron and apply the perceptron learning rule.
        \item Distinguish single-layer from multilayer networks and explain why hidden layers matter.
        \item Explain the forward pass, loss functions, and training loop conceptually.
        \item Articulate what backpropagation does and how the chain rule makes it efficient.
    \end{itemize}
\end{frame}

\section{Biological Inspiration}

\begin{frame}{Why Neural Networks?}
    Neural networks represent a shift from symbolic AI to learning from data.

    \vspace{0.2cm}
    \textbf{Key motivations:}
    \begin{itemize}
        \item Brain-inspired: learning through example, adaptability.
        \item Universal approximators: can learn any function (given enough capacity).
        \item Powerful for high-dimensional data (images, audio, text).
        \item Foundation for deep learning and modern AI.
    \end{itemize}

    \vspace{0.2cm}
    \textit{Conceptual goal:} Understand how neurons work together to learn patterns.
\end{frame}

\begin{frame}{The Biological Neuron}
    \textbf{Structure:}
    \begin{itemize}
        \item \textbf{Dendrites:} Receive signals from other neurons.
        \item \textbf{Cell body (soma):} Integrates incoming signals.
        \item \textbf{Axon:} Transmits signal to other neurons.
        \item \textbf{Synapses:} Junctions where signals pass; strength varies through learning.
    \end{itemize}

    \vspace{0.2cm}
    \textbf{Key insight:} Learning means adjusting synaptic weights through experience.

    \vspace{0.2cm}
    \textbf{Artificial simplification:}
    \begin{itemize}
        \item Inputs $\leftrightarrow$ dendrites
        \item Weights $\leftrightarrow$ synaptic strengths
        \item Activation function $\leftrightarrow$ firing threshold
        \item Output $\leftrightarrow$ axon signal
    \end{itemize}
\end{frame}

\begin{frame}{From Neurons to Networks}
    A single neuron is too simple to be useful alone.

    \vspace{0.2cm}
    \textbf{Scale in biology:}
    \begin{itemize}
        \item Human brain: $\sim$86 billion neurons.
        \item Each neuron connects to thousands of others.
        \item Intelligence emerges from massive, structured connectivity.
    \end{itemize}

    \vspace{0.2cm}
    \textbf{Scale in artificial networks:}
    \begin{itemize}
        \item Layers of artificial neurons organised in parallel.
        \item Early networks: tens of neurons.
        \item Modern LLMs: billions of parameters.
    \end{itemize}

    \vspace{0.2cm}
    \textit{The architecture and learning algorithm matter as much as size.}
\end{frame}

\section{The Perceptron}

\begin{frame}{The Perceptron}
    \textbf{Simplest artificial neuron (Rosenblatt, 1958):}
    \[
        \text{output} = \begin{cases}
            1 & \text{if } \displaystyle\sum_i w_i x_i + b \geq 0 \\
            0 & \text{otherwise}
        \end{cases}
    \]

    \vspace{0.2cm}
    \textbf{Components:}
    \begin{itemize}
        \item \textbf{Inputs} $x_1, \ldots, x_n$: feature values.
        \item \textbf{Weights} $w_1, \ldots, w_n$: importance of each feature (learned).
        \item \textbf{Bias} $b$: shifts the decision threshold.
        \item \textbf{Activation:} step function (outputs 0 or 1).
    \end{itemize}

    \vspace{0.2cm}
    \textbf{Geometric view:} The perceptron defines a \emph{hyperplane} that separates two classes.
\end{frame}

\begin{frame}{Perceptron Learning Rule}
    \textbf{How does a perceptron learn?}

    \vspace{0.2cm}
    \textbf{Algorithm:}
    \begin{enumerate}
        \item Initialise weights $w_i$ and bias $b$ to small random values.
        \item For each training example $(x, y)$:
        \begin{itemize}
            \item Compute prediction: $\hat{y} = \text{step}\!\left(\sum_i w_i x_i + b\right)$
            \item If $\hat{y} \neq y$ (wrong):
            \[
                w_i \;\gets\; w_i + \alpha\,(y - \hat{y})\,x_i, \quad b \gets b + \alpha\,(y-\hat{y})
            \]
            where $\alpha > 0$ is the \emph{learning rate}.
        \end{itemize}
        \item Repeat until convergence (or max epochs).
    \end{enumerate}

    \vspace{0.2cm}
    \textbf{Key idea:} Nudge weights toward correct classification on each mistake.
\end{frame}

\begin{frame}{Perceptron Limitations: The XOR Problem}
    \textbf{Fatal limitation:} Perceptrons can only learn \textbf{linearly separable} patterns.

    \vspace{0.2cm}
    \textbf{XOR truth table:}
    \small
    \begin{tabularx}{0.6\textwidth}{c c c}
        \toprule
        $x_1$ & $x_2$ & $x_1 \oplus x_2$ \\
        \midrule
        0 & 0 & 0 \\
        0 & 1 & 1 \\
        1 & 0 & 1 \\
        1 & 1 & 0 \\
        \bottomrule
    \end{tabularx}

    \normalsize
    \vspace{0.2cm}
    No single straight line can separate the two classes.

    \vspace{0.2cm}
    \textbf{Impact (1969):} Minsky \& Papert's analysis caused the first ``AI winter'' for neural nets.

    \vspace{0.2cm}
    \textbf{Solution:} Stack multiple perceptrons into a \emph{multilayer network}.
\end{frame}

\section{Multilayer Neural Networks}

\begin{frame}{Multilayer Network Architecture}
    \textbf{Three types of layers:}
    \begin{itemize}
        \item \textbf{Input layer:} Raw feature values (no computation).
        \item \textbf{Hidden layer(s):} Learn intermediate representations.
        \item \textbf{Output layer:} Produces final prediction.
    \end{itemize}

    \vspace{0.2cm}
    \textbf{Why hidden layers?}
    \begin{itemize}
        \item Stacking linear functions is still linear --- no gain.
        \item Non-linear activation functions break linearity at each layer.
        \item Hidden layers learn progressively abstract features.
    \end{itemize}

    \vspace{0.2cm}
    \textbf{Depth vs. width:}
    \begin{itemize}
        \item \emph{Deep} network: many layers (complex hierarchical features).
        \item \emph{Wide} network: many units per layer (more capacity per layer).
    \end{itemize}
\end{frame}

\begin{frame}{Activation Functions}
    \textbf{Purpose:} Introduce non-linearity between layers.

    \vspace{0.2cm}
    \textbf{Common choices:}
    \begin{itemize}
        \item \textbf{Step:} $\mathbf{1}[z \geq 0]$ --- original perceptron; not differentiable.
        \item \textbf{Sigmoid:} $\sigma(z) = \dfrac{1}{1+e^{-z}}$ --- smooth, range $(0,1)$.
        \item \textbf{Tanh:} $\tanh(z)$ --- symmetric around 0, range $(-1, 1)$.
        \item \textbf{ReLU:} $\max(0, z)$ --- simple, fast; default in modern networks.
    \end{itemize}

    \vspace{0.2cm}
    \textbf{Why differentiability matters:} Gradient-based training requires computing $\frac{d}{dz}[\text{activation}]$.

    \vspace{0.2cm}
    \textit{ReLU's simplicity and avoidance of vanishing gradients make it the go-to choice.}
\end{frame}

\begin{frame}{Forward Pass}
    \textbf{Goal:} Compute the network's output given an input.

    \vspace{0.2cm}
    For layer $l$:
    \[
        z^{(l)} = W^{(l)}\, a^{(l-1)} + b^{(l)}, \qquad a^{(l)} = \sigma\!\left(z^{(l)}\right)
    \]
    where $W^{(l)}, b^{(l)}$ are the weight matrix and bias, $a^{(l)}$ is the layer output.

    \vspace{0.2cm}
    \textbf{Loss function:} Measures how wrong the prediction is.
    \begin{itemize}
        \item \textbf{Regression} (MSE): $L = \dfrac{1}{m}\displaystyle\sum_i (y_i - \hat{y}_i)^2$
        \item \textbf{Classification} (cross-entropy):
        $L = -\dfrac{1}{m}\displaystyle\sum_i\bigl[y_i \log \hat{y}_i + (1-y_i)\log(1-\hat{y}_i)\bigr]$
    \end{itemize}

    \vspace{0.2cm}
    \textbf{Goal of training:} Minimise $L$ over all training examples.
\end{frame}

\section{Backpropagation}

\begin{frame}{The Training Problem}
    \textbf{Challenge:} We have weights across many layers. How do we know which weights to change and by how much?

    \vspace{0.3cm}
    \textbf{Gradient descent} updates weights in the direction that reduces loss:
    \[
        w \;\gets\; w - \alpha\,\frac{\partial L}{\partial w}
    \]
    where $\alpha$ is the \emph{learning rate}.

    \vspace{0.3cm}
    \textbf{Problem:} Computing $\frac{\partial L}{\partial w}$ for every weight in a deep network is expensive if done naively.

    \vspace{0.2cm}
    \textbf{Solution:} \textbf{Backpropagation} --- efficiently compute all gradients in a single backward pass using the chain rule.
\end{frame}

\begin{frame}{Backpropagation: Conceptual Overview}
    \textbf{Four steps:}
    \begin{enumerate}
        \item \textbf{Forward pass:} Compute predictions and loss $L$.
        \item \textbf{Output gradient:} Compute $\frac{\partial L}{\partial \hat{y}}$ at the output layer.
        \item \textbf{Backward pass:} Propagate error signals layer-by-layer \emph{backward} through the network.
        \item \textbf{Weight update:} Adjust every weight $w \gets w - \alpha\,\frac{\partial L}{\partial w}$.
    \end{enumerate}

    \vspace{0.2cm}
    \textbf{Key property:} Each layer only needs the gradient from the layer \emph{above} it --- no redundant computation.

    \vspace{0.2cm}
    \textit{Backprop makes training deep networks computationally feasible.}
\end{frame}

\begin{frame}{Backpropagation: The Chain Rule}
    \textbf{Core mechanism:} The chain rule of calculus.

    \vspace{0.2cm}
    Example --- gradient w.r.t.\ weights in layer 1 of a 2-layer network:
    \[
        \frac{\partial L}{\partial W^{(1)}} =
        \underbrace{\frac{\partial L}{\partial \hat{y}}}_{\text{output error}}
        \cdot
        \underbrace{\frac{\partial \hat{y}}{\partial z^{(2)}}}_{\sigma'}
        \cdot
        \underbrace{\frac{\partial z^{(2)}}{\partial a^{(1)}}}_{{W^{(2)}}}
        \cdot
        \underbrace{\frac{\partial a^{(1)}}{\partial z^{(1)}}}_{\sigma'}
        \cdot
        \underbrace{\frac{\partial z^{(1)}}{\partial W^{(1)}}}_{a^{(0)}}
    \]

    \vspace{0.2cm}
    \textbf{Intuition:}
    \begin{itemize}
        \item Start with the error at the output.
        \item Multiply through each layer's local gradient.
        \item Each layer ``receives'' how much it contributed to the error.
    \end{itemize}
\end{frame}

\begin{frame}{Training Loop}
    \textbf{Standard training procedure:}
    \begin{enumerate}
        \item \textbf{Initialise} weights randomly (small values).
        \item For each \emph{epoch} (pass over training data):
        \begin{enumerate}
            \item \textbf{Forward:} Compute $\hat{y}$ and loss $L$.
            \item \textbf{Backward:} Compute all gradients via backprop.
            \item \textbf{Update:} $W \gets W - \alpha\,\nabla_W L$.
        \end{enumerate}
        \item Monitor loss on a \emph{validation set}; stop when it no longer improves.
    \end{enumerate}

    \vspace{0.2cm}
    \textbf{Key hyperparameters:}
    \begin{itemize}
        \item Learning rate $\alpha$: too large $\to$ diverge; too small $\to$ slow.
        \item Number of layers and units per layer.
        \item Batch size (how many examples per update).
        \item Regularisation strength.
    \end{itemize}
\end{frame}

\begin{frame}{Overfitting and Generalisation}
    \textbf{Overfitting:} The model memorises training data and performs poorly on new data.

    \vspace{0.2cm}
    \textbf{Signs:} Training loss decreases, but validation loss increases.

    \vspace{0.2cm}
    \textbf{Common remedies:}
    \begin{itemize}
        \item \textbf{More data:} Usually the most effective fix.
        \item \textbf{Dropout:} Randomly zero out units during training (forces redundancy).
        \item \textbf{L2 regularisation:} Penalise large weights: $L_{\text{reg}} = L + \lambda \|W\|^2$.
        \item \textbf{Early stopping:} Stop training when validation loss stops improving.
        \item \textbf{Batch normalisation:} Normalise layer inputs; stabilises training.
    \end{itemize}
\end{frame}

\begin{frame}{Strengths and Challenges}
    \textbf{Strengths:}
    \begin{itemize}
        \item \textbf{Universal approximation:} A network with one hidden layer can approximate any continuous function (given enough units).
        \item Learn complex non-linear patterns automatically.
        \item End-to-end learning: raw input to prediction without hand-crafted features.
    \end{itemize}

    \vspace{0.2cm}
    \textbf{Challenges:}
    \begin{itemize}
        \item \textbf{Local minima \& saddle points:} Gradient descent is not guaranteed to find the global optimum.
        \item \textbf{Vanishing/exploding gradients:} Gradients can shrink or grow exponentially in deep nets.
        \item \textbf{Interpretability:} Hard to explain individual decisions (black box).
        \item \textbf{Data and compute:} Large networks need large datasets and significant compute.
    \end{itemize}
\end{frame}

\begin{frame}{Summary}
    \textbf{What we covered:}
    \begin{itemize}
        \item Biological inspiration: neurons, synapses, and learning as weight adjustment.
        \item Perceptron: linear classifier, learning rule, XOR limitation.
        \item Multilayer networks: hidden layers, activation functions, forward pass.
        \item Backpropagation: chain rule, efficient gradient computation, training loop.
        \item Practical concerns: overfitting, regularisation, hyperparameter tuning.
    \end{itemize}

    \vspace{0.2cm}
    \textbf{Key takeaway:} Neural networks learn by adjusting millions of weights to minimise a loss function --- backpropagation makes this tractable.

    \vspace{0.2cm}
    \textbf{Further study:}
    \begin{itemize}
        \item Convolutional Neural Networks (CNNs) for images.
        \item Recurrent Neural Networks (RNNs) for sequences.
        \item Transformers and attention mechanisms for NLP.
    \end{itemize}
\end{frame}

\end{document}
